\subsection{Finite support repair preserving a lower norm class}

\begin{statement}{A.100}{Proved}
{Relative ray-class repair in one existing finite norm chain}
Fix the odd-prime data of A.95 with $\ell^t\mid M$, and
choose compatible positive roots of the current rational
coefficient $a$.
For $0\le s<t$ put
\[
 E=\Q(a^{1/\ell^t}),\qquad
 K=\Q(a^{1/\ell^s})\subset E,\qquad
 L=E(\mu_{\ell^t}).
 \tag{A.100.1}
\]
Here $K=\Q$ when $s=0$.
For the nested fields of A.77 these are the same actual
subfields after the explicit rational rescaling of roots
in (A.86.1).
Let $\beta\in E^\times$ have rational norm $b$, and put
$\beta_s=\mathrm N_{E/K}\beta$.
Fix a precision $\ell^r$, a finite rational-prime set $S$
containing two, every prime ramified in $L/\Q$, and the
rational primes below the support of $\beta_s$.
Let $\mathfrak m$ be a finite ray modulus of $E$ supported
above $S\setminus\{2\}$.
Retain the prescribed cyclotomic Frobenius $\rho$ of A.95.

There is $\beta'\in E^\times$ with the same exact rational
norm $b$, the same valuations above $S$, and
$\beta'/\beta\equiv1\pmod{\mathfrak m}$, such that all
valuations outside $S$ and finitely many permitted primes
$Q$ are divisible by $\ell^r$.
Every $q\in Q$ splits completely in $L$ and the current
$B$-radical field, has the prescribed cyclotomic Frobenius,
and may avoid any fixed finite set and lie above any fixed bound.
At each such $q$, the new valuation vector has sum zero
\emph{separately above each prime of $K$ over $q$}.
Moreover there is an actual unit $\eta\in\mathcal O_K^\times$
such that
\[
                  \mathrm N_{E/K}\beta'
                                =\beta_s\eta^{\ell^r}.
 \tag{A.100.2}
\]
Consequently, for every field $K_i$ in the same chain
below $K$, the exact identity
\[
 \frac{\mathrm N_{E/K_i}\beta'}{\mathrm N_{E/K_i}\beta}
                           =(\mathrm N_{K/K_i}\eta)^{\ell^r}
 \tag{A.100.3}
\]
preserves the lower norm class modulo $\ell^r$.
This repairs the support of an already existing finite
chain.  It does not supply a new norm preimage of an
arbitrarily fixed lower representative.
\end{statement}
\begin{proof}
We prove relative ideal-class generation with its exact
subgroups, then arrange that the resulting relative norm
unit is a global $\ell^r$th power.

\emph{A finite collection of unit-detecting places.}
The group
$\mathcal O_K^\times/(\mathcal O_K^\times)^{\ell^r}$
is finite by the unit theorem.
For odd $\ell^r$, the local-power theorem
\cite[VIII, Theorem 1.4]{MilneCFT}, printed page 229,
says that an element which is an $\ell^r$th power at
all but finitely many places is a global $\ell^r$th power.
If a unit is a power in $K^\times$, its root is also
a unit, by taking valuations.
Thus every nonzero class in the displayed finite unit
quotient is detected at some finite place outside any
prescribed finite set.
Choose one such place for each nonzero class, avoiding
the primes above two and $\ell$, the original set $S$,
and every rational prime below the support of $\beta$.
Write $T$ for these finitely many places of $K$.
The actual localization map
\[
 \mathcal O_K^\times/(\mathcal O_K^\times)^{\ell^r}
       \longrightarrow
       \prod_{v\in T}K_v^\times/(K_v^\times)^{\ell^r}
                   \quad\text{is injective}.
 \tag{A.100.4}
\]
Real and complex places are unnecessary for odd powers.
If the unit quotient is trivial, take $T$ empty.

Enlarge $S$ temporarily by the rational primes below
$T$ and call the result $S^*$.
The element $\beta$ has valuation zero at every prime
above $S^*\setminus S$, by the choice just made.
Enlarge $\mathfrak m$ to a ray modulus $\mathfrak m^*$
containing every prime of $E$ above each $v\in T$ to
at least the first power.
It still has no prime above two.
If $z\in E^\times$ is a unit above $T$ and is one
modulo these primes, then
$\mathrm N_{E/K}z\in1+\mathfrak p_v$ at every $v\in T$.
The group $1+\mathfrak p_v$ consists of $\ell^r$th powers,
because the residue characteristic is prime to $\ell$
and Hensel lifting applies to $X^{\ell^r}-u$ at one.
In particular, if $\mathrm N_{E/K}z$ is a global unit,
these ray conditions and (A.100.4) make it a global
$\ell^r$th-power unit.

\emph{The relative Artin groups and their generators.}
Let $H_{\mathfrak m^*}/E$ be the primary ray class field,
put $F=L\prod_{\sigma\in G_\Q}\sigma(H_{\mathfrak m^*})$,
and define
\[
 G=\operatorname{Gal}(F/\Q),\quad
 J=\operatorname{Gal}(F/K),\quad
 H=\operatorname{Gal}(F/E),\quad
 V=\operatorname{Gal}(F/L).
 \tag{A.100.5}
\]
Restriction gives $H$ onto
$C=\operatorname{Cl}_{\mathfrak m^*}(E)\{\ell\}$.
Write $U_t=(\Z/\ell^t\Z)^\times$.
On the root labels, fixing $a^{1/\ell^s}$ means that
the translation is divisible by $\ell^s$.
Thus the exact subquotients are
\[
 J/V=(\ell^s\Z/\ell^t\Z)\rtimes U_t,
                    \qquad H/V=U_t.
 \tag{A.100.6}
\]
The translation group and $U_t$ are cyclic.
A generator of $U_t$ has $u-1$ invertible modulo $\ell$,
so the cyclic calculation of A.92 gives
\[
 H^2(J/V,\Q_\ell/\Z_\ell)=0,
            \qquad (J/V)^{\rm ab}=U_t.
 \tag{A.100.7}
\]
The map $H/V\to(J/V)^{\rm ab}$ is the identity on
this $U_t$.

Define $C_{\rm rel,split}$ to be the subgroup of $C$
generated by differences of prime ideals of $E$ above
the \emph{same} prime of $K$, at rational primes
splitting completely in $L$.
We claim that every relative ideal-norm-zero divisor
supported outside $S^*$ has class in this subgroup,
and that its generators can satisfy all the boundary
and cyclotomic conditions of A.95.

Here are the character maps proving the claim.
For a character $\psi:C\to\Q_\ell/\Z_\ell$, pull it
back to $H$ and set $\lambda=\psi|_V$.
At a rational prime split in $L$, choose Frobenius
$\sigma\in V$ in $F$.
For the prime of $K$ below this choice, the primes
of $E$ above it are indexed by $J/H$, and their
character values are
$\lambda(j^{-1}\sigma j)$ for $jH\in J/H$.
Every $\sigma\in V$ is supplied by Chebotarev outside
any finite set.
Consequently annihilation of these same-prime
differences is exactly $J$-invariance of $\lambda$.
The five-term sequence and (A.100.7) extend that
$\lambda$ to a character of $J$.
Its difference from $\psi$ on $H$ factors through
$H/V=U_t$ and extends through $J\to U_t$.
We therefore obtain a character $\chi$ of $J$ with
$\chi|_H=\psi$ literally.

For any prime $v$ of $K$ outside $S^*$, choose its
Frobenius $\tau\in J$.
Its cycles on $J/H$ have the relative residue
degrees $f_i$ of the primes of $E$ above $v$.
Exactly as in (A.92.9), their Artin character values
are $f_i\chi(\tau)$.
Every relative ideal-norm-zero vector satisfies
$\sum f_i t_i=0$, so $\psi$ kills it.
Finite character separation proves the claimed
relative generation.

The restricted auxiliary primes give the same subgroup.
Indeed, $F/\Q$ remains unramified at two, because the
ray modulus excludes two.
The field intersection (A.95.7) and its automorphism
construction apply unchanged.
The permitted Frobenius elements in $V$ form a
nonempty coset $\sigma_0V_0$ whose quotient $V/V_0$
has trivial $G$-action, hence trivial $J$-action.
If $\psi$ annihilates the permitted same-prime
differences, the functions
$\lambda(j^{-1}vj)-\lambda(v)$ vanish on $V_0$.
They define a homomorphism
$J\to\operatorname{Hom}(V/V_0,\Q_\ell/\Z_\ell)$,
as in (A.95.10).
It kills $V$ and $H$, and (A.100.7) makes it zero.
Thus $\lambda$ is $J$-invariant and $\psi$ annihilates
all of $C_{\rm rel,split}$.
This proves the claim with the prescribed boundary
splitting, cyclotomic Frobenius, and finite exclusions.

\emph{The ray-principal correction and its lower norm.}
Write $\operatorname{div}(\beta)=D_{S^*}+D_{\rm aux}$.
At every prime of $K$ outside $S^*$ the relative
ideal norm of its part of $D_{\rm aux}$ is zero:
it equals the valuation of $\beta_s$, whose support
was included in $S$.
There is no contribution at primes in $S^*$.
Hence $\mathrm N_{E/K}D_{\rm aux}=0$ literally.
The proved generation gives an integral divisor
$D_Q$ supported at finitely many permitted fresh
rational primes, with
\[
 \mathrm N_{E/K}D_Q=0,\qquad
                 [D_Q]=[D_{\rm aux}]\quad\text{in }C.
 \tag{A.100.8}
\]
Its norm is zero separately above each prime of $K$;
this is why the same-prime generators were retained.
As in A.99, equality in the primary ray class group
gives an integer $d$ prime to $\ell$ and a ray-principal
element $u_0$ with
\[
 \operatorname{div}(u_0)=d(D_{\rm aux}-D_Q),
                         \qquad u_0\equiv1\pmod{\mathfrak m^*}.
 \tag{A.100.9}
\]
Its relative norm has zero divisor, so it is an
actual unit of $K$.
The detector places make it $\eta_0^{\ell^r}$ for
some $\eta_0\in\mathcal O_K^\times$.
Its rational norm is $\pm1$.
Replace $u_0$ by $u=u_0^2$ and $d$ by $2d$.
Then
\[
 \mathrm N_{E/\Q}u=1,
       \qquad \mathrm N_{E/K}u=(\eta_0^2)^{\ell^r},
       \qquad \operatorname{div}(u)=d(D_{\rm aux}-D_Q).
 \tag{A.100.10}
\]
Choose $t_0d\equiv1\pmod{\ell^r}$ and put
$\beta'=\beta u^{-t_0}$.
Its divisor is
$D_{S^*}+(1-t_0d)D_{\rm aux}+t_0dD_Q$,
so the same exact support and ray calculations as
in (A.99.8) apply.
At the temporarily added detecting primes,
$D_{S^*}$ is zero.  Thus the final support conclusion
uses the original $S$, as stated.
The relative norm is
$\beta_s(\eta_0^{-2t_0})^{\ell^r}$, which is (A.100.2)
with $\eta=\eta_0^{-2t_0}$.
Its norms to the lower fields give (A.100.3).
This proves the finite-precision assertion using one
fixed finite chain and its actual norm maps.
\end{proof}

\noindent The exact certificate
\path{code/norm_relative_ray_certificate.py}, with retained output
\path{code/norm_relative_ray_certificate_output.txt}, checks the
relative holomorph and root-stabilizer maps, the residue-degree
Artin formula, the same-prime lower norm identities, and the
literal power exponents in (A.100.2)--(A.100.3).
The finite unit-detecting places are obtained from the cited
odd-power local-global theorem in the proof.
